\chapter{Prestazioni di base}
\label{cap:23-prestazioni-base}

Questo capitolo riporta le prestazioni della piattaforma sul carico di
riferimento, attraverso sei release consecutive, sul profilo hardware fissato
descritto nel \cref{cap:22-metodo-sperimentale}.

\section{La serie}

Tutte le misure che seguono sono mediane di tre esecuzioni sul profilo
\code{azure-d8s-v5+d2s-v5-amd64-native-loadtest-v1}: control plane compilato
nativamente per AMD64 su VM \code{Standard\_D8s\_v5}, generatore di carico k6 su
VM separata \code{Standard\_D2s\_v5}, scenario \code{autoscaling-cycle-k8s}.

\begin{table}[htbp]
\centering
\caption{Serie di prestazioni per release. Mediane di tre esecuzioni sullo
stesso profilo.}
\label{tab:serie}
\small
\begin{adjustbox}{max width=\textwidth}
\begin{tabular}{@{}lrrrr@{}}
\toprule
\textbf{Versione} & \textbf{Throughput (req/s)} & \textbf{Tasso errore} &
\textbf{p95 (ms)} & \textbf{Repliche di picco} \\
\midrule
0.17.0  & 241{,}90 & 0{,}000 & 6{,}50  & 3 \\
0.17.1  & 269{,}07 & 0{,}000 & 5{,}77  & 3 \\
v0.18.1 & 239{,}13 & 0{,}000 & 6{,}61  & 4 \\
v0.18.2 & 249{,}34 & 0{,}000 & 5{,}13  & 3 \\
v0.18.3 & 232{,}41 & 0{,}000 & 5{,}13  & 4 \\
v0.19.0 & \textbf{300{,}48} & 0{,}000 & \textbf{11{,}26} & 5 \\
\bottomrule
\end{tabular}
\end{adjustbox}
\end{table}

\begin{figure}[htbp]
  \centering
  \resizebox{0.95\textwidth}{!}{% Release performance series, from docs/performance/history.md (medians of 3 runs
% on profile azure-d8s-v5+d2s-v5-amd64-native-loadtest-v1).
\documentclass[border=8pt]{standalone}
\usepackage{nanofaas-figs}
\usepackage{pgfplots}
\pgfplotsset{compat=1.18}
\begin{document}
\begin{tikzpicture}
\begin{axis}[
    width=13cm, height=7cm,
    axis y line*=left, axis x line*=bottom,
    xlabel={release}, ylabel={throughput (req/s)},
    ylabel style={nfblue}, y tick label style={nfblue},
    symbolic x coords={0.17.0,0.17.1,v0.18.1,v0.18.2,v0.18.3,v0.19.0},
    xtick=data, ymin=0, ymax=340,
    font=\sffamily\small, tick label style={font=\sffamily\footnotesize},
    legend style={draw=none, font=\sffamily\footnotesize, at={(0.02,0.98)}, anchor=north west},
    ybar, bar width=16pt,
]
  \addplot[fill=nfbluebg, draw=nfblue, line width=0.8pt] coordinates {
    (0.17.0,241.904) (0.17.1,269.069) (v0.18.1,239.130)
    (v0.18.2,249.341) (v0.18.3,232.407) (v0.19.0,300.483)};
  \addlegendentry{throughput}
\end{axis}
\begin{axis}[
    width=13cm, height=7cm,
    axis y line*=right, axis x line=none,
    ylabel={p95 (ms)}, ylabel style={nfamber}, y tick label style={nfamber},
    symbolic x coords={0.17.0,0.17.1,v0.18.1,v0.18.2,v0.18.3,v0.19.0},
    ymin=0, ymax=13, font=\sffamily\small,
    tick label style={font=\sffamily\footnotesize},
    legend style={draw=none, font=\sffamily\footnotesize, at={(0.02,0.86)}, anchor=north west},
]
  \addplot[nfamber, line width=1.3pt, mark=*, mark size=2.2pt] coordinates {
    (0.17.0,6.496) (0.17.1,5.772) (v0.18.1,6.607)
    (v0.18.2,5.125) (v0.18.3,5.131) (v0.19.0,11.256)};
  \addlegendentry{p95}
\end{axis}
\end{tikzpicture}
\end{document}
}
  \caption{Throughput (barre, asse sinistro) e p95 (linea, asse destro) per
  release. Dati da \code{docs/performance/history.md}.}
  \label{fig:serie}
\end{figure}

\section{Che cosa la serie mostra}

\subsection{Il tasso di errore \`e esattamente zero}

Su tutte e sei le release, su tre esecuzioni ciascuna. Non ``prossimo a zero'':
zero. \`E il gate rigido del flusso di rilascio, e la sua tenuta \`e il
risultato pi\`u solido del capitolo --- l'unico che non dipende dalla
sensibilit\`a del metodo, perch\'e distingue fra zero e non-zero anzich\'e fra
due valori vicini.

Va ricordato (\cref{cap:21-osservabilita}) che questo contatore esclude le
risposte non-2xx deliberate delle funzioni, grazie al meccanismo della busta.
Un gate a zero sarebbe inapplicabile senza quella disambiguazione.

\subsection{Il throughput oscilla in una banda, poi sale}

Le prime cinque release stanno fra 232 e 269 richieste al secondo: una banda del
$\pm 7\%$ attorno a 250. Con tre esecuzioni per punto, \emph{questa banda non
\`e distinguibile dal rumore}. Nessuna delle differenze fra 0.17.0, v0.18.1,
v0.18.2 e v0.18.3 pu\`o essere attribuita a una modifica del codice sulla base
di questi dati.

La v0.19.0 esce dalla banda: 300{,}5 req/s, ossia $+29\%$ sulla release
precedente e $+12\%$ sul massimo precedente. Questa differenza \`e maggiore di
qualunque oscillazione osservata prima, e supera la soglia del gate.

\subsection{Il p95 della v0.19.0 \`e il dato da spiegare}

Nella stessa release il p95 passa da 5{,}13\,ms a 11{,}26\,ms: pi\`u che
raddoppiato, e ben oltre la soglia di regressione del 15\%. Il numero di
repliche di picco sale da 4 a 5.

L'interpretazione che i dati sostengono \`e che la v0.19.0 abbia spostato il
sistema in un punto di lavoro diverso, non che sia peggiorata. Un throughput di
300 req/s su cinque repliche significa un carico per replica confrontabile con
quello precedente; il p95 pi\`u alto \`e coerente con un sistema che opera pi\`u
vicino alla propria capacit\`a. Il record di dettaglio della v0.19.0 lo conferma:

\begin{center}
\small
\begin{adjustbox}{max width=\textwidth}
\begin{tabular}{@{}lr@{}}
\toprule
\textbf{Grandezza} & \textbf{v0.19.0} \\
\midrule
p50 & 4{,}43\,ms \\
p95 & 11{,}26\,ms \\
p99 & 21{,}05\,ms \\
attesa in coda (media) & 0{,}82\,ms \\
cold start osservati & 5 \\
CPU di picco del control plane & 0{,}664 core \\
heap di picco del control plane & 199{,}5\,MB \\
\bottomrule
\end{tabular}
\end{adjustbox}
\end{center}

Il p50 di 4{,}43\,ms \`e in linea con le release precedenti: il tipico non \`e
peggiorato, si \`e allungata la coda della distribuzione. L'attesa in coda media
\`e sotto il millisecondo, quindi la coda non \`e la causa.

\begin{damisurare}
L'attribuzione del salto di p95 della v0.19.0 non \`e stata condotta. I dati
disponibili escludono l'attesa in coda come causa dominante e sono coerenti con
un punto di lavoro pi\`u vicino alla saturazione, ma non lo stabiliscono. Un
confronto a parit\`a di throughput offerto fra v0.18.3 e v0.19.0 lo
deciderebbe.
\todomis{p95 di v0.18.3 e v0.19.0 a parit\`a di carico offerto}
\end{damisurare}

\subsection{L'osservabilit\`a del control plane}

Due grandezze del record riguardano il control plane e non le funzioni: CPU di
picco 0{,}664 core e heap di picco 199{,}5\,MB. Sono la risposta quantitativa
alla domanda del \cref{cap:01-introduzione} su quanto costi un control plane
minimale: meno di un core e circa 200\,MB di heap per sostenere 300 richieste al
secondo.

Va sottolineato che si tratta di \emph{heap}, non di memoria residente. Il
\cref{cap:24-footprint} mostra che per un processo Spring su JVM la memoria
residente \`e diverse volte l'heap, perch\'e la maggior parte non \`e dati.

\section{Gli SLO indicativi}

Il progetto dichiara obiettivi di lavoro, esplicitamente non contrattuali:

\begin{center}
\small
\begin{adjustbox}{max width=\textwidth}
\begin{tabular}{@{}lp{7.6cm}@{}}
\toprule
\textbf{Obiettivo} & \textbf{Stato} \\
\midrule
Tasso di errore 0\% & Rispettato in tutte e sei le release; \`e un gate rigido. \\
p95 $\leq$ 10\,ms & Rispettato fino a v0.18.3; \textbf{violato} in v0.19.0
(11{,}26\,ms). \\
Throughput $\geq$ 200 req/s & Rispettato in tutte le release. \\
Cold start & Nessun budget fisso; si traccia la distribuzione. \\
\bottomrule
\end{tabular}
\end{adjustbox}
\end{center}

La scelta di non fissare un budget per il cold start ma di tracciarne la
distribuzione \`e coerente con l'analisi del \cref{cap:02-stato-dell-arte}: il
cold start su \nf{} \`e dominato dal tempo di creazione del pod da parte di
Kubernetes e dall'avvio del runtime, entrambi fuori dal controllo del control
plane. Un budget su una grandezza che non si governa \`e un obiettivo che si
manca senza poter reagire.

\begin{quote}
\textbf{Nota sulla documentazione.} Il documento \code{docs/slo.md} riporta la
serie fino a v0.18.3 e non include la v0.19.0, che \`e presente in
\code{docs/performance/history.md} e nel record di release. La discrepanza \`e
segnalata qui perch\'e la violazione dell'obiettivo di p95 \`e visibile solo nella
seconda fonte. \`E un esempio del genere di deriva che colpisce la
documentazione mantenuta a mano, gi\`a osservato per \code{openapi.yaml} nel
\cref{cap:05-contratti-dati}.
\end{quote}

\section{Limiti di questa serie}

\begin{enumerate}[leftmargin=*]
  \item \textbf{Tre esecuzioni per punto.} Non consentono di stimare una
  varianza. Le differenze inferiori al 10\% non sono interpretabili.

  \item \textbf{Infrastruttura condivisa.} Due VM Azure dello stesso tipo non
  hanno necessariamente le stesse prestazioni. Parte dell'oscillazione osservata
  fra le prime cinque release potrebbe essere questa.

  \item \textbf{Un solo scenario.} Tutta la serie usa
  \code{autoscaling-cycle-k8s}. Non esistono serie comparabili per altri profili
  di carico o altre famiglie di funzioni.

  \item \textbf{Il confronto \`e a coppie consecutive.} Il gate confronta con il
  record precedente, quindi una deriva lenta che resti sotto soglia a ogni passo
  non viene rilevata. La serie completa nella \cref{tab:serie} \`e il rimedio
  manuale, ma non \`e automatizzato.
\end{enumerate}

Il quarto punto merita enfasi perch\'e \`e un difetto strutturale del metodo, non
dei dati: un gate che confronta solo con l'immediato precedente ammette una
regressione monotona del 9\% per release indefinitamente.

\begin{damisurare}
Non esiste una verifica automatica contro un riferimento assoluto o contro la
migliore release storica. La \cref{tab:serie} rende la deriva visibile a chi la
legge, e nient'altro.
\todomis{gate contro un riferimento storico oltre al precedente immediato}
\end{damisurare}
